\documentclass[11pt]{article}

\usepackage[T1]{fontenc}
\usepackage[utf8]{inputenc}
\usepackage[english]{babel}
\usepackage[margin=1in]{geometry}
\usepackage{lmodern}
\usepackage{microtype}
\usepackage{mathtools}
\usepackage{amssymb}
\usepackage{amsthm}
\usepackage{graphicx}
\usepackage{booktabs}
\usepackage{array}
\usepackage{caption}
\usepackage{enumitem}
\usepackage{fancyhdr}
\usepackage[numbers,sort&compress]{natbib}
\usepackage{orcidlink}
\usepackage{titling}
\usepackage{doi}
\usepackage{xurl}
\usepackage{hyperref}

\hypersetup{
  pdftitle={Strict Theory Extension on a Lawful Continuous Cantor Shell},
  pdfauthor={Ioannis Tsiokos},
  hidelinks
}
\urlstyle{same}
\captionsetup{
  font=small,
  labelfont=bf,
  labelsep=period,
  margin=1em
}
\captionsetup[table]{skip=6pt,position=top}
\captionsetup[figure]{skip=6pt}
\setlist[enumerate]{leftmargin=*, itemsep=2pt, topsep=4pt}
\setlist[itemize]{leftmargin=*, itemsep=2pt, topsep=4pt}
\setlength{\droptitle}{-3.2em}
\pretitle{\begin{center}\LARGE\bfseries}
\posttitle{\par\end{center}\vspace{0.5em}}
\preauthor{\begin{center}\normalsize}
\postauthor{\par\end{center}\vspace{-0.4em}}
\predate{\begin{center}\small}
\postdate{\par\end{center}\vspace{0.6em}}
\setlength{\emergencystretch}{2em}
\clubpenalty=10000
\widowpenalty=10000
\displaywidowpenalty=10000
\interfootnotelinepenalty=10000
\allowdisplaybreaks[2]
\setlength{\bibsep}{2pt plus 0.3ex}

\pagestyle{fancy}
\fancyhf{}
\fancyhead[L]{\small\itshape Canonical Hybrid Cantor Object}
\fancyhead[R]{\small\thepage}
\renewcommand{\headrulewidth}{0.4pt}
\renewcommand{\footrulewidth}{0pt}

\fancypagestyle{firstpage}{\fancyhf{}
  \fancyfoot[C]{\scriptsize Automorph Inc., Wilmington, DE, USA \quad \texttt{ioannis@automorph.io} \quad \mbox{DOI: \href{https://doi.org/10.5281/zenodo.19189353}{10.5281/zenodo.19189353}} \quad \textcopyright\ 2026 \quad CC-BY 4.0 \quad Preprint --- not peer reviewed.}
  \renewcommand{\headrulewidth}{0pt}
  \renewcommand{\footrulewidth}{0pt}
}

\usepackage{titlesec}
\titlespacing*{\section}{0pt}{1.8ex plus 0.8ex minus 0.4ex}{1.0ex plus 0.3ex}
\titlespacing*{\subsection}{0pt}{1.4ex plus 0.6ex minus 0.3ex}{0.7ex plus 0.2ex}

\newcommand{\classid}[1]{\texttt{\small\hyphenchar\font=`\_ #1}}

\newtheorem{theorem}{Theorem}[section]
\newtheorem{proposition}[theorem]{Proposition}
\newtheorem{lemma}[theorem]{Lemma}
\newtheorem{corollary}[theorem]{Corollary}
\theoremstyle{definition}
\newtheorem{definition}[theorem]{Definition}
\newtheorem{remark}[theorem]{Remark}

\newcommand{\R}{\mathbb{R}}
\newcommand{\N}{\mathbb{N}}
\newcommand{\Tzero}{T_0}
\newcommand{\Tone}{T_1}
\newcommand{\Hybrid}{\mathfrak{H}}
 
\begin{document}

\title{Strict Theory Extension on a\\Lawful Continuous Cantor Shell}
\author{Ioannis Tsiokos\,\orcidlink{0009-0009-7659-5964}}
\date{\today}

\maketitle
\thispagestyle{firstpage}

\begin{abstract}
We study a class of continuous Cantor substrates generated by six interacting mechanisms---operator rewrite, admissibility gating, protocol/timescale adaptation, lens selection, packaging/completion, and budget dynamics---and isolate an audited shell on which the resulting full-loop dynamics are lawful. On this shell we define a base cocycle-pressure theory~$T_0$ and a completion-based extension~$T_1$ built from an evolve--forget--reinstantiate packaging endomap. We prove four theoremlets on this shell-stable class: (i)~a continuous full-loop lawfulness theoremlet, (ii)~closure of a selector-weighted cocycle pressure object, (iii)~a strict theory-extension theoremlet showing that~$T_1$ is not definable from~$T_0$, via saturation, material $P4 \leftarrow P5$ forcing, and macro-admissibility obstruction, and (iv)~a thermodynamic consequence theoremlet in the form of a nontrivial conditional pressure disintegration over packaging fibers. The contribution is therefore a theorem of theory depth rather than broader family coverage: the packaged object carries object structure not present in the cocycle object even on a fixed shell-stable class. We make no broader-class, shell-general, or non-SFT breadth claim beyond this audited shell, but the resulting hybrid Cantor object yields a closed and thermodynamically meaningful audited-shell theorem package.
\end{abstract}

\smallskip

\noindent\textbf{Keywords:} Cantor dynamics; thermodynamic formalism; theory extension; pressure closure; conditional disintegration; six birds theory.

\medskip

\paragraph{Contribution summary.}
The paper contributes a closed audited-shell theorem package for a canonical hybrid Cantor object: a continuous full-loop lawfulness theoremlet, a cocycle pressure closure theoremlet, a strict theory extension theoremlet, and a conditional pressure disintegration theoremlet.
 \section{Introduction}

Classical Cantor and fractal constructions are usually built from fixed symbolic rules, fixed iterated function systems, or fixed graph-directed dynamics \cite{Hutchinson1981Fractals,Moran1946AdditiveFunctions,MauldinUrbanskiGDMS}. In such settings the main mathematical objects are already given in advance: one starts with a rule system and then studies the sets, measures, and thermodynamic quantities generated by that rule system in the standard fractal-geometry and thermodynamic-formalism sense \cite{FalconerFractalGeometry,Bowen1979Quasicircles,Ruelle1978ThermodynamicFormalism,WaltersErgodic1982}. The present paper starts from a different problem. We ask whether a \emph{lawful evolving Cantor substrate} can itself support a mathematically stable object theory, a thermodynamic theory, and a genuine notion of theory extension without collapsing into a finite toy model or dissolving into unconstrained simulation.

The difficulty is not merely dynamical. It is object-theoretic. If one only keeps the cocycle-level description of the substrate, one obtains a clean pressure-based theory but risks missing object structure that appears only after completion/packaging. If one only keeps packaged fixed points, one loses the thermodynamic side. The task is therefore to identify a theorem object suitable for the audited-shell package. In the present paper, that object is a \emph{canonical hybrid Cantor object} defined on an audited shell-stable class of lawful continuous substrates: a base cocycle-pressure theory \(T_0\), together with a completion-based extension \(T_1\) obtained from an evolve--forget--reinstantiate packaging endomap. The resulting object is neither a raw trajectory nor a static completion summary; it is a coupled theorem object in which cocycle structure and packaged fixed-point structure coexist.

The substrate itself is generated by six interacting mechanisms:
\[
\begin{aligned}
&\text{P1: operator rewrite},\qquad
\text{P2: admissibility gating},\qquad
\text{P3: protocol/timescale adaptation},\\
&\text{P4: lens selection},\qquad
\text{P5: packaging/completion},\qquad
\text{P6: budget/ledger dynamics}.
\end{aligned}
\]
These mechanisms define the substrate construction \cite{TsiokosFoundations}. The mathematical results below are stated in terms of the resulting cocycle, completion map, packaged strata, and shell-stable class.

The motivating contrast is with adaptive regimes in which the effective maps may depend on location, time, or local context \cite{OliveiraVarandasLocalIFS,BarnsleyDemkoEltonGeronimo1988,FalconerRandomFractals1986,RempeGillenUrbanski2016NonAutonomousCIFS}. The present paper does not pursue that broader theory. It isolates one audited shell on which the coupled object theory closes.

The paper develops a closed theorem package on an audited shell-stable class of continuous full-loop Cantor substrates. Its novelty lies in \emph{theory depth}, not broader family coverage. In particular, the main result is not a broad new theorem class beyond all symbolic or finite-state regimes. Rather, it is that a cocycle-pressure theory closes on a lawful shell, that packaging completion defines additional object structure on that same shell, and that this extended object has a nontrivial thermodynamic consequence.

More precisely, the paper closes the following four theoremlets on the audited shell:
\begin{enumerate}
    \item a \emph{continuous full-loop lawfulness theoremlet},
    \item a \emph{cocycle pressure closure theoremlet},
    \item a \emph{strict theory extension theoremlet}, and
    \item a \emph{conditional pressure disintegration theoremlet}.
\end{enumerate}
The first two establish a lawful and thermodynamically meaningful base theory \(T_0\). The third proves that packaging completion produces a strict extension \(T_1\) of that base theory, via saturation, material \(P4 \leftarrow P5\) forcing, non-definability from \(T_0\), and macro-admissibility obstruction. The fourth shows that the thermodynamic consequence of this extension is not direct stratumwise root separation, but a nontrivial conditional pressure disintegration over packaging fibers.

\paragraph{Contributions.}
The paper makes four main contributions.

\begin{enumerate}
    \item \textbf{A lawful continuous Cantor substrate on an audited shell.}
    We identify a shell-stable class of continuous full-loop Cantor substrates on which all six mechanisms remain causally active and the dynamics do not collapse into a trivial frozen regime.

    \item \textbf{A closed cocycle-pressure theory on that shell.}
    On this class we define a selector-weighted cocycle pressure object and close its pressure theoremlet on the audited shell.

    \item \textbf{A strict theory-extension theorem for the canonical hybrid object.}
    We define a completion-based extension of the cocycle theory and show that, on the audited shell, the resulting hybrid object is a strict theory extension of the cocycle theory rather than a mere re-description of it.

    \item \textbf{A thermodynamic consequence in the form of conditional pressure disintegration.}
    We show that the correct thermodynamic consequence of the extension is a nontrivial weighted package-conditioned pressure gap, rather than direct stratumwise root separation.
\end{enumerate}

Equally important is what the paper \emph{does not} claim. We do not claim a broader-class theorem beyond the audited shell. We do not claim a shell-general theorem. We do not claim an external or non-SFT breadth theorem beyond what is actually closed here. We do not claim a direct stratumwise root-separation theorem. And we do not claim that the packaging route has already produced a broader theorem class than the cocycle route. The paper is therefore intentionally conservative in scope and explicit about object-level closure.

The conceptual lesson is that progress in this branch is not best measured by larger family coverage. It is measured by the appearance of a deeper theorem object. The cocycle object \(T_0\) is thermodynamically closed, but too coarse to support closed macro-dynamics for the packaged future. The completion-based object \(T_1\) supplies the missing object structure. The mathematical result is therefore a theorem of theory depth: the packaged object is strictly richer than the cocycle object even on a fixed shell-stable class.

\paragraph{Organization of the paper.}
Section~\ref{sec:preliminaries} fixes notation and the audited-shell setting. Section~\ref{sec:canonical-object} defines the canonical hybrid object and distinguishes it from the raw simulation trajectory. Section~\ref{sec:lawfulness} proves the lawfulness theoremlet. Section~\ref{sec:cocycle-pressure} proves cocycle pressure closure on the audited shell. Section~\ref{sec:strict-extension} proves that the completion-based theory is a strict extension of the cocycle theory. Section~\ref{sec:conditional-disintegration} proves the conditional pressure disintegration consequence. Section~\ref{sec:related-work} positions the result relative to both Cantor/fractal dynamics and prior substrate-generation work. Section~\ref{sec:discussion} records the scope restrictions, non-claims, and future directions. Appendix~\ref{sec:appendix-supporting-evidence} summarizes the supporting audited-shell evidence that accompanies, but does not replace, the theorem package.
 \section{Preliminaries}\label{sec:preliminaries}

This section fixes the mathematical setting and notation used throughout the paper. The central point is that the theorem object of the paper is not the raw simulation trajectory of the continuous substrate, but a canonical object built from two coupled levels of description: a cocycle-pressure theory \(\Tzero\) and a completion-based extension \(\Tone\). Standard background references for the surrounding concepts include classical IFS/fractal geometry \cite{Hutchinson1981Fractals,FalconerFractalGeometry,MauldinUrbanskiGDMS,Moran1946AdditiveFunctions}, thermodynamic formalism \cite{Ruelle1978ThermodynamicFormalism,WaltersErgodic1982,BarreiraNonadditiveTF}, and subadditive or almost-additive extensions of that formalism \cite{CaoFengHuang2008SubAdditive,Mummert2006AlmostAdditive}.

\subsection{Continuous substrate states and the audited shell}

A \emph{continuous substrate state} is written
\[
x=(K,\ell,\mathfrak{p},\tau,b,\eta)\in \mathcal{X},
\]
where:
\begin{itemize}
    \item \(K\in\mathcal{K}\subset \R^{d\times d}\) is the current normalized kernel/operator state,
    \item \(\ell\in\mathcal{L}\) is the currently active lens state,
    \item \(\mathfrak{p}\in\mathcal{P}\) is the currently active packaging state,
    \item \(\tau\in\mathcal{T}\subset \R_{>0}\) is the current timescale/protocol parameter,
    \item \(b\in\mathcal{B}\subset \R\) is the budget/ledger variable,
    \item \(\eta\in\mathcal{H}\) denotes the remaining auxiliary selector and hysteresis state.
\end{itemize}
The full update law is written
\[
F:\mathcal{X}\to\mathcal{X},\qquad x\mapsto F(x),
\]
and is generated by the six interacting mechanisms listed in the introduction:
operator rewrite (\(P1\)), admissibility gating (\(P2\)), protocol/timescale adaptation (\(P3\)), lens selection (\(P4\)), packaging/completion (\(P5\)), and budget/ledger dynamics (\(P6\)).

We work on a distinguished audited shell-stable class
\[
\mathcal{S}_{\mathrm{aud}}\subset \mathcal{X},
\]
which is the class denoted in the internal package by
\[
\texttt{continuous\_full\_loop\_lawful\_kernel\_class\_shell\allowbreak\_stable}.
\]
All closed theoremlets in this paper are statements on \(\mathcal{S}_{\mathrm{aud}}\). No broader-class theorem, shell-general theorem, or non-SFT breadth claim is made beyond this audited shell.

\subsection{The base theory \texorpdfstring{$\Tzero$}{T0}}

The base theory \(\Tzero\) is the cocycle-pressure theory carried by the audited shell. We write
\[
\pi_0:\mathcal{S}_{\mathrm{aud}}\to \mathcal{O}_0
\]
for the \(\Tzero\)-object map, where \(\mathcal{O}_0\) denotes the cocycle-level object space. Intuitively, \(\pi_0(x)\) records the object visible at the cocycle/lens level before completion-based packaging is adjoined.

The thermodynamic object attached to \(\Tzero\) is a cocycle-pressure quantity. For the chosen observable family \(a_n(s;x)\), depending on depth/time \(n\), parameter \(s\), and shell state \(x\), the associated pressure is written
\[
P_{\Tzero}(s)
\]
once the closure theoremlet has established that the corresponding pressure object is well-defined on the audited shell. At the level of notation we suppress the shell-state dependence whenever the shell-stable theorem package has already fixed it.

\subsection{The completion extension \texorpdfstring{$\Tone$}{T1}}

The extended theory \(\Tone\) adjoins packaging completion to the cocycle theory. For a given lens \(\ell\) and timescale \(\tau\), the packaging endomap is
\[
E_{\tau,\ell}(\mu):=U_{\ell}\!\bigl(Q_{\ell}(\mu K^{\tau})\bigr),
\]
where:
\begin{itemize}
    \item \(\mu\) is a state/distribution object on which the completion acts,
    \item \(Q_{\ell}\) is the lens-dependent forgetting/coarse-graining map,
    \item \(U_{\ell}\) is the corresponding reinstatement/prototype lift.
\end{itemize}
The fixed points of \(E_{\tau,\ell}\) are the packaged objects of the completion theory. We write
\[
\mathrm{Fix}(E_{\tau,\ell})
\]
for the fixed-point set and
\[
\pi_1:\mathcal{S}_{\mathrm{aud}}\to \mathcal{O}_1
\]
for the \(\Tone\)-object map, where \(\mathcal{O}_1\) is the packaged-object space.

A \emph{packaged stratum} is a persistent element or equivalence class in the packaged fixed-point structure associated with the current shell state. We write
\[
\mathcal{F}(x)
\]
for the finite family of persistent packaged strata determined by \(x\in \mathcal{S}_{\mathrm{aud}}\).

\subsection{The canonical hybrid object}

The theorem object of the paper is not \(\pi_0(x)\) alone and not \(\pi_1(x)\) alone. It is the canonical hybrid object obtained by combining the cocycle-level object map, the packaging-completion endomap, and the packaged strata. We write it schematically as
\[
\mathfrak{H}(x)=\bigl(\pi_0(x),\,\pi_1(x),\,\mathcal{F}(x)\bigr).
\]
The role of \(\mathfrak{H}\) is twofold:
\begin{enumerate}
    \item it keeps track of the closed cocycle-pressure theory \(\Tzero\), and
    \item it records the genuinely richer packaged-object structure used to define the extension \(\Tone\).
\end{enumerate}
In particular, strict theory extension will be formulated as a statement that the packaged-object map does not factor through the cocycle-level object map on the audited shell.

\subsection{Pressure and conditional disintegration objects}

The cocycle pressure \(P_{\Tzero}(s)\) is the base thermodynamic quantity. On the completion side, the paper does \emph{not} use direct stratumwise root separation as its thermodynamic consequence. Instead, the consequence object is a weighted package-conditioned pressure gap. If \(\mathcal{F}(x)\) is the family of packaged strata and \(w_{\Sigma}(x)\) are the associated package weights, then the consequence quantity is written
\[
\Delta_{\Tzero\to \Tone}(s)
:=
P_{\Tzero}(s)-\sum_{\Sigma\in\mathcal{F}(x)} w_{\Sigma}(x)\,P_{\Sigma}(s),
\]
where \(P_{\Sigma}(s)\) denotes the package-conditioned pressure profile attached to the stratum \(\Sigma\). The conditional-disintegration theoremlet states that this weighted gap is nontrivial on the audited shell.

The role of \(\Delta_{\Tzero\to \Tone}(s)\) is that it measures the thermodynamic insufficiency of the cocycle-level theory for the packaged future. In the paper’s final consequence theoremlet, the relevant thermodynamic statement is therefore a \emph{conditional disintegration statement}, not a direct theorem of independent stratumwise root separation \cite{Rokhlin1949FundamentalIdeas}.

\subsection{Definability, saturation, and obstruction}

Two further pieces of notation are used repeatedly.

First, \(\Tzero\)-\emph{definability} means definability from the cocycle-level object/lens structure alone. At the formal level this is expressed through factorization of the packaged-object map through \(\pi_0\). Failure of such factorization on the audited shell is the non-definability signal behind the strict theory extension theoremlet.

Second, the completion dynamics may \emph{saturate}: repeated application of \(E_{\tau,\ell}\) ceases to produce new packaged objects at the current theory depth. Material \(P4\leftarrow P5\) forcing then changes the lens and reveals new packaged strata in the extended theory.

Finally, we use \emph{macro-admissibility obstruction} to mean that the cocycle-level theory is too coarse to support closed macro-dynamics on the packaged future. This obstruction is one of the certificates that \(\Tone\) is a strict theory extension of \(\Tzero\).

\subsection{Scope convention}

Unless explicitly stated otherwise, every theorem statement in the paper is made on the audited shell \(\mathcal{S}_{\mathrm{aud}}\) and for the canonical hybrid object \(\mathfrak{H}\). We do not claim:
\begin{itemize}
    \item a broader-class theorem beyond the audited shell,
    \item a shell-general theorem,
    \item a non-SFT or external breadth theorem beyond what is actually closed,
    \item a direct stratumwise root-separation theorem,
    \item or a packaging-induced broader theorem class claim.
\end{itemize}
The paper is therefore about a closed theorem package on a fixed shell-stable class, with novelty coming from \emph{theory depth} rather than family breadth.
 \section{The canonical hybrid object}\label{sec:canonical-object}

The purpose of this section is to identify the mathematical object that the paper is actually about. The key point is that the theorem object is neither the raw trajectory of the continuous substrate nor the cocycle-level object alone nor the completion structure alone. It is a \emph{canonical hybrid object} in which a closed cocycle-pressure theory and a completion-based packaged-object theory coexist on the same audited shell.

\subsection{Why the raw trajectory is not the theorem object}

The continuous substrate evolves by the full six-mechanism update law
\[
x\mapsto F(x), \qquad x\in \mathcal{S}_{\mathrm{aud}}.
\]
This trajectory is the source of the mathematical structures we study, but it is not itself the theorem object. A theorem stated directly about raw trajectories would be too tied to the simulation-level presentation and would not distinguish between:
\begin{itemize}
    \item the cocycle-level object visible through the base theory,
    \item the packaged fixed-point structure visible only after completion,
    \item and the extension relation between those two levels.
\end{itemize}
What the paper needs is an object that remembers both the closed thermodynamic structure already present at the cocycle level and the new packaged strata created by completion.

\subsection{The base theory \texorpdfstring{$\Tzero$}{T0} as an object map}

The base theory \(\Tzero\) associates to each shell state \(x\in \mathcal{S}_{\mathrm{aud}}\) an object
\[
\pi_0(x)\in \mathcal{O}_0,
\]
where \(\mathcal{O}_0\) is the cocycle-level object space introduced in Section~\ref{sec:preliminaries}. The map
\[
\pi_0:\mathcal{S}_{\mathrm{aud}}\to \mathcal{O}_0
\]
forgets the completion/fixed-point structure and retains only the cocycle-level description relevant to the closed pressure theoremlet. The resulting object theory is already nontrivial: on the audited shell, the associated selector-weighted cocycle pressure object closes and gives a mathematically stable thermodynamic description.

However, the base theory is deliberately coarse. It does not, by itself, distinguish packaged strata created by the completion dynamics. In the language of the later strict-extension theoremlet, \(\Tzero\) is the base theory whose object map is too coarse to express the full packaged future.

\subsection{The completion theory \texorpdfstring{$\Tone$}{T1} as a packaged-object map}

The extended theory \(\Tone\) augments the cocycle-level object theory by the packaging endomap \cite{TsiokosCreateStone,TsiokosFoundations}
\[
E_{\tau,\ell}(\mu)=U_{\ell}\!\bigl(Q_{\ell}(\mu K^{\tau})\bigr).
\]
For each shell state \(x\), this completion mechanism yields a packaged-object map
\[
\pi_1(x)\in \mathcal{O}_1,
\qquad
\pi_1:\mathcal{S}_{\mathrm{aud}}\to \mathcal{O}_1,
\]
where \(\mathcal{O}_1\) is the packaged-object space. Concretely, \(\pi_1(x)\) records the persistent fixed-point strata and their organization after completion, saturation, and material \(P4\leftarrow P5\) refinement.

Thus \(\Tzero\) and \(\Tone\) do not differ merely by carrying two different observables on the same object. They define different object maps:
\[
\pi_0:\mathcal{S}_{\mathrm{aud}}\to \mathcal{O}_0,
\qquad
\pi_1:\mathcal{S}_{\mathrm{aud}}\to \mathcal{O}_1.
\]
The strict-extension theoremlet later proves that \(\pi_1\) is not definable from \(\pi_0\) on the audited shell.

\subsection{The canonical hybrid object}

The theorem object of the paper is the pair consisting of the base object and the completion object on the same shell state, together with the packaged fixed-point family. We write this schematically as
\[
\Hybrid(x)
=
\bigl(\pi_0(x),\,\pi_1(x),\,\mathcal{F}(x)\bigr),
\qquad x\in \mathcal{S}_{\mathrm{aud}}.
\]
Here \(\mathcal{F}(x)\) denotes the persistent packaged strata associated with \(x\). The notation suppresses the explicit dependence on the completion endomap \(E_{\tau,\ell}\), but that endomap is part of the structure that defines \(\pi_1\) and \(\mathcal{F}(x)\).

This object is canonical for the present theorem package for three reasons.

\begin{enumerate}
    \item It retains the base cocycle-pressure object needed for the pressure closure theoremlet.
    \item It retains the packaged fixed-point object needed for the completion and extension theoremlets.
    \item It records both at the same shell state, making it possible to ask whether the packaged object is definable from the cocycle object.
\end{enumerate}

In particular, the hybrid object is the smallest natural object that can support all four closed theoremlets at once.

\subsection{Definability as factorization}

The correct notion of definability for the extension problem is factorization of object maps \cite{LindMarcus1995SymbolicDynamics,Walters1986RelativePressure}. If the packaged-object map were definable from the cocycle-level object map on the audited shell, then there would exist a map
\[
\phi:\mathcal{O}_0\to \mathcal{O}_1
\]
such that
\[
\pi_1=\phi\circ\pi_0
\qquad\text{on }\mathcal{S}_{\mathrm{aud}}.
\]
In that case the completion theory would add no genuinely new object structure: the packaged object would be recoverable from the base cocycle object.

The point of the strict-extension theoremlet is that this factorization fails on the audited shell. Equivalently, there exist shell states with the same \(\Tzero\)-object but different packaged-object strata in \(\Tone\). The later theorem does not use broader family coverage to certify novelty. It uses failure of factorization.

\subsection{Saturation, forcing, and packaged strata}

The completion theory is not static. The packaged strata arise through a sequence of completion operations that may saturate at a fixed theory depth. Saturation means that repeated completion at the current lens/timescale no longer produces new packaged objects. Material \(P4\leftarrow P5\) forcing then changes the lens and reveals additional packaged strata at the extended theory depth.

This is why the hybrid object is richer than the cocycle object. The cocycle side supplies a closed thermodynamic base, but the packaged side carries fixed-point strata whose appearance depends on saturation and forcing. The hybrid object therefore has both a thermodynamic face and an extension-theoretic face.

\subsection{Role in the theorem package}

The closed theoremlets of the paper live at different levels of the same object hierarchy:
\begin{itemize}
    \item the lawfulness theoremlet is a theorem about the audited shell carrying the full update law,
    \item the cocycle pressure closure theoremlet is a theorem about the \(\Tzero\)-component \(\pi_0(x)\),
    \item the strict theory extension theoremlet is a theorem about the relation between \(\pi_0(x)\) and \(\pi_1(x)\),
    \item the conditional disintegration theoremlet is a theorem about thermodynamic quantities on the hybrid object \(\Hybrid(x)\).
\end{itemize}
Thus the canonical hybrid object is not an optional repackaging of the paper’s content. It is the only object on which the full theorem package can be stated coherently.

\subsection{Scope}

The object hierarchy above is fixed on the audited shell \(\mathcal{S}_{\mathrm{aud}}\). We do not claim:
\begin{itemize}
    \item that the same hybrid object theorem package holds on a broader class,
    \item that the packaged-object map yields a broader theorem class than the cocycle theory,
    \item that direct stratumwise root separation is the correct thermodynamic object,
    \item or that any shell-general or non-SFT breadth theorem has been proved here.
\end{itemize}
The paper’s novelty lies instead in the fact that, even on a fixed shell-stable class, the packaged object is strictly richer than the cocycle object and supports an additional thermodynamic consequence.
 \section{Lawfulness on the audited shell}\label{sec:lawfulness}

We now state the first closed theoremlet in the paper: the continuous full-loop lawfulness theoremlet on the audited shell. This result is the base structural theorem in the package. It does not yet concern strict theory extension or thermodynamic consequence. Its role is to show that the shell-stable class supports a well-defined, nondegenerate, six-mechanism update law whose dynamics do not collapse immediately into a trivial frozen regime.

\subsection{Lawful shell data}

Recall from Section~\ref{sec:preliminaries} that a substrate state is
\[
x=(K,\ell,\mathfrak{p},\tau,b,\eta)\in \mathcal{S}_{\mathrm{aud}},
\]
with continuous operator/kernel state \(K\), active lens \(\ell\), active packaging \(\mathfrak{p}\), timescale \(\tau\), budget variable \(b\), and auxiliary selector/hysteresis state \(\eta\). The update law
\[
F:\mathcal{S}_{\mathrm{aud}}\to \mathcal{S}_{\mathrm{aud}}
\]
is generated by the six interacting mechanisms \(P1\)--\(P6\), viewed here as a deterministic hybrid update with switching and hysteresis on the audited shell \cite{GoebelSanfeliceTeel2012Hybrid}. The shell-stable class is the class denoted internally by
\[
\texttt{continuous\_full\_loop\_lawful\_kernel\_class\_shell\allowbreak\_stable}.
\]

For manuscript purposes, the hypotheses that matter here are the following:
\begin{enumerate}
    \item the operator state remains in a controlled shell of kernel states,
    \item all six mechanisms remain causally active on that shell,
    \item lens and packaging selectors remain nondegenerate under hysteresis,
    \item the budget and timescale variables remain in a lawful range,
    \item the resulting dynamics exhibit persistent variation and do not undergo fast collapse.
\end{enumerate}
These hypotheses are exactly the shell-level assumptions frozen in the internal theorem package; no broader class is claimed in this section.

\begin{theorem}[Continuous full-loop lawfulness on the audited shell]\label{thm:lawfulness}
Let \(\mathcal{S}_{\mathrm{aud}}\) be the audited shell-stable class of continuous substrate states described above. Assume that the full update law \(F\) satisfies the shell-stable hypotheses just listed: well-posed continuous operator dynamics, full six-mechanism causal closure, hysteresis-stable lens and packaging selection, lawful budget/timescale evolution, and absence of fast collapse on the shell. Then the following hold on \(\mathcal{S}_{\mathrm{aud}}\).

\begin{enumerate}
    \item The update law is well-defined and remains inside the shell.
    \item Lens competition and packaging competition remain nondegenerate.
    \item Budget and timescale dynamics remain active and lawful.
    \item All six mechanisms \(P1\)--\(P6\) remain causally active in the resulting dynamics.
    \item The shell supports a nontrivial lawful exploratory regime rather than immediate collapse into a degenerate frozen state.
\end{enumerate}

In particular, \(\mathcal{S}_{\mathrm{aud}}\) supports a continuous full-loop lawful regime.
\end{theorem}

\begin{proof}[Proof sketch]
The proof is a synthesis of the closed shell arguments from the internal dependency package. The point is not to reprove every implementation-level estimate inside the paper, but to expose the mathematical structure of the closure.

\smallskip
\noindent\emph{Step 1: well-posed state space.}
The state variables \((K,\ell,\mathfrak{p},\tau,b,\eta)\) define a well-posed update state on the audited shell. The kernel state evolves in a controlled continuous class, while the selector, budget, and timescale variables live in corresponding bounded shell ranges. This is the well-posedness ingredient behind the state-space lemma of the internal package.

\smallskip
\noindent\emph{Step 2: refresh and selector consistency.}
The refresh/invalidation mechanism is consistent on the shell, so lens and packaging candidates can be recomputed without ambiguity. The hysteretic selectors therefore define stable lens and packaging dynamics rather than arbitrary switching. In particular, lens and packaging competition remain meaningful rather than collapsing into undefined or pathological selection.

\smallskip
\noindent\emph{Step 3: six-mechanism closure.}
The shell hypotheses ensure that each of the six mechanisms enters the update law essentially:
\(P1\) changes the operator state,
\(P2\) changes admissibility/support,
\(P3\) changes timescale/protocol state,
\(P4\) changes the active lens,
\(P5\) changes the active packaging and feeds into other update components,
and \(P6\) changes the budget/ledger state.
Thus the lawful loop is not a partial mechanism loop disguised as a full one.

\smallskip
\noindent\emph{Step 4: forward shell preservation.}
The shell assumptions imply that one step of the update law keeps the dynamics inside the lawful shell. In particular, kernel variation, selector margins, budget range, and timescale range remain in the admissible region. This is the shell-preservation ingredient later reused by the theorem package.

\smallskip
\noindent\emph{Step 5: noncollapse.}
Because selector competition, budget dynamics, and timescale adaptation remain active inside the shell, the resulting dynamics do not collapse into a trivial frozen regime. The audited shell therefore carries a nontrivial lawful exploratory regime rather than a transient simulation artifact.

Combining these points proves the theoremlet.
\end{proof}

\begin{figure}[t]
\centering
\small
\begin{tabular}{@{}l r@{}}
\toprule
\multicolumn{2}{@{}l}{\textbf{Continuous substrate regime on the audited shell}} \\
\midrule
Base witness kernel dimension        & 16 \\
Shell witness kernel dimension       & 20 \\
Base witness lens switches           & 249 \\
Base witness packaging switches      & 249 \\
Base witness budget range            & $[2.714,\;12.0]$ \\
Base witness $\tau$ range            & $[0.600,\;0.982]$ \\
Shell witness budget range           & $[2.739,\;12.0]$ \\
Shell witness $\tau$ range           & $[0.600,\;0.861]$ \\
Shell exits on both witnesses        & 0 \\
All six primitives active            & yes \\
\bottomrule
\end{tabular}
\caption{Continuous full-loop regime on the audited shell.
The two frozen witnesses remain inside the shell-stable class, exhibit persistent
variation, nondegenerate lens and packaging competition, active
budget/timescale dynamics, and no shell exits.
This figure supports only the continuous full-loop lawfulness
theoremlet on the audited shell.}
\label{fig:full-loop-regime}
\end{figure}
 \begin{figure}[t]
\centering
\small
\begin{tabular}{@{}l c c l@{}}
\toprule
\multicolumn{4}{@{}l}{\textbf{Primitive knockout comparison on the audited shell}} \\
\midrule
Removed & Lens sw. & Pkg.\ sw. & Main degradation \\
\midrule
none & 249 & 249 & reference lawful loop \\
P1   & 249 & 249 & $\tau$ switching collapse \\
P2   & 249 & 249 & stabilization change and damping \\
P3   &   0 &   0 & selector and $\tau$ collapse \\
P4   &   0 & 249 & lens collapse \\
P5   & 249 &   0 & packaging, budget, and $\tau$ degradation \\
P6   & 249 & 249 & budget collapse \\
\bottomrule
\end{tabular}
\caption{Leave-one-primitive-out comparison on the audited shell.
The full loop remains lawful, while removal of any one primitive
materially degrades the regime.
In particular, removing $P4$ collapses lens switching,
removing $P5$ collapses packaging competition and narrows the budget range,
and removing $P6$ collapses budget dynamics.
This figure supports only the six-primitive closure component
of the lawfulness theoremlet.}
\label{fig:primitive-knockout-closure}
\end{figure}
 
The importance of Theorem~\ref{thm:lawfulness} is that it supplies the base dynamical regime on which all later objects are built. The cocycle-pressure theoremlet, the strict theory-extension theoremlet, and the conditional disintegration theoremlet all presuppose that the shell carries a lawful full loop. Without Theorem~\ref{thm:lawfulness}, the later theoremlets would sit on a moving or ill-defined substrate.

We emphasize again the scope of the result. The theoremlet is a statement on the audited shell only. It does not assert a shell-general lawfulness theorem, a broader-class theorem, or any external/non-SFT breadth statement. Its role is to close the base lawful regime on which the rest of the paper is built.
 \section{Cocycle pressure closure}\label{sec:cocycle-pressure}

This section states the second closed theoremlet in the paper: cocycle pressure closure on the audited shell. The theoremlet belongs entirely to the base theory \(\Tzero\). Its purpose is to show that the shell-stable lawful regime from Section~\ref{sec:lawfulness} already carries a well-defined thermodynamic object before completion-based packaging is adjoined.

\subsection{The pressure route}

The pressure route fixed in the internal theorem package is the \emph{Fekete-sup cocycle route} \cite{Fekete1923DistributionRoots} with the selected observable family given by the \emph{selector-weighted operator growth observable}. The point of this choice is to keep the pressure object tied to the lawful shell as a genuine six-mechanism object rather than to a frozen symbolic proxy.

For each shell state \(x\in\mathcal{S}_{\mathrm{aud}}\), parameter \(s\), and horizon \(n\), let
\[
a_n(s;x)
\]
denote the selected finite-horizon cocycle observable. We do not need its full implementation-level formula in the paper; what matters here is that it is built from the cocycle/operator evolution together with the active selector structure on the shell. We then define the shell-level envelope
\[
\Phi_n(s):=\sup_{x\in\mathcal{S}_{\mathrm{aud}}} a_n(s;x).
\]
The base-theory pressure candidate is
\[
P_{\Tzero}(s):=\lim_{n\to\infty}\frac{1}{n}\log \Phi_n(s),
\]
whenever the limit exists. The cocycle-pressure closure theoremlet states that this pressure object is well-defined on the audited shell and has the regularity needed for the later thermodynamic package. The broader context is the standard subadditive and nonadditive thermodynamic-formalism literature \cite{BarreiraNonadditiveTF,CaoFengHuang2008SubAdditive,Mummert2006AlmostAdditive}, together with the familiar role of subadditive limits in ergodic theory \cite{Kingman1968Subadditive}.

\subsection{Why the pressure object belongs to \texorpdfstring{$\Tzero$}{T0}}

At this stage of the paper the relevant object is still the cocycle-level object map
\[
\pi_0:\mathcal{S}_{\mathrm{aud}}\to \mathcal{O}_0.
\]
The completion map \(E_{\tau,\ell}\), packaged fixed-point strata, and strict extension relation belong to later sections. The role of the present theoremlet is therefore deliberately narrower: it closes the thermodynamic object attached to the base cocycle theory before the packaging-completion theory is introduced as a richer object level.

In particular, this section does \emph{not} claim:
\begin{itemize}
    \item a pressure theorem beyond the audited shell,
    \item a shell-general pressure theorem,
    \item a theorem about the packaged-object space \(\mathcal{O}_1\),
    \item or any direct stratumwise thermodynamic distinction theorem.
\end{itemize}
Its scope is exactly the shell-stable base theory \(\Tzero\).

\begin{theorem}[Cocycle pressure closure on the audited shell]\label{thm:cocycle-pressure}
Let \(\mathcal{S}_{\mathrm{aud}}\) be the audited shell-stable class from Section~\ref{sec:lawfulness}, and let \(a_n(s;x)\) be the selector-weighted operator growth observable associated with the base cocycle theory \(\Tzero\). Define
\[
\Phi_n(s):=\sup_{x\in\mathcal{S}_{\mathrm{aud}}} a_n(s;x).
\]
Assume the shell-stable hypotheses of the lawfulness theoremlet, together with the shell-uniform cocycle bounds and bounded-error comparison estimates required by the Fekete-sup route. Then the following hold on the audited shell.

\begin{enumerate}
    \item The pressure object
    \[
    P_{\Tzero}(s)=\lim_{n\to\infty}\frac{1}{n}\log \Phi_n(s)
    \]
    exists on the relevant parameter window.
    \item \(P_{\Tzero}(s)\) is finite and nondegenerate on that window.
    \item \(P_{\Tzero}(s)\) is monotone in \(s\).
    \item \(P_{\Tzero}(s)\) has enough regularity to support a later root-based thermodynamic theorem.
\end{enumerate}

Therefore the cocycle-pressure object of the base theory \(\Tzero\) is closed on the audited shell.
\end{theorem}

\begin{proof}[Proof sketch]
The proof follows the closed pressure-dependency package for the shell-stable class.

\smallskip
\noindent\emph{Step 1: cocycle observable well-posedness.}
Because the full-loop substrate is lawful on \(\mathcal{S}_{\mathrm{aud}}\), the selected cocycle observable \(a_n(s;x)\) is well-defined for every shell state \(x\), parameter \(s\), and time horizon \(n\). This uses the shell-level control established in Section~\ref{sec:lawfulness}.

\smallskip
\noindent\emph{Step 2: shell-uniform control.}
The audited shell provides uniform finite-horizon control on the selected observable. In particular, the cocycle growth observable remains bounded in the sense needed for the pressure route, and does not escape the shell through selector collapse, budget collapse, or loss of timescale control.

\smallskip
\noindent\emph{Step 3: bounded-error comparison and the Fekete step.}
The shell-uniform estimates imply a bounded-error comparison of Fekete type for the logarithmic envelope \(\log \Phi_n(s)\). This is the key almost-subadditive input. It yields existence of the pressure object through the chosen Fekete-sup route.

\smallskip
\noindent\emph{Step 4: finiteness and nondegeneracy.}
The same shell estimates rule out both blow-up and trivial collapse on the parameter window relevant to the theorem package. Hence the pressure object is finite and nondegenerate there.

\smallskip
\noindent\emph{Step 5: monotonicity and root-supporting regularity.}
Monotonicity in \(s\) follows from the monotonic dependence built into the selected observable family. Combined with the shell-uniform control, this yields the regularity needed for the later root-based thermodynamic consequence statements. From the thermodynamic-formalism side, this is the point where the pressure object should be compared conceptually with standard transfer-operator and operator-growth settings \cite{Ruelle1978ThermodynamicFormalism,Baladi2000TransferOperators}; likewise, the cocycle-growth viewpoint sits naturally beside the classical matrix-product context of \cite{FurstenbergKesten1960RandomMatrices,Oseledets1968MET}, even though the present theoremlet is proved by the audited-shell Fekete route.

These five steps prove the theoremlet.
\end{proof}

\begin{table}[t]
\caption{Support values for the closed cocycle-pressure theoremlet on the audited shell.
The Fekete-sup route remains finite and nondegenerate on both witnesses,
with positive pressure proxies and shell-uniform cocycle control.
This table supports only the cocycle pressure closure theoremlet
on the audited shell.}
\label{tbl:pressure-closure-support}
\centering
\small
\begin{tabular}{@{}l c c c c c c@{}}
\toprule
Witness & Dim & Shell exits & Fekete gap & Growth bound & Lens margin & Pkg.\ margin \\
\midrule
Base  & 16 & 0 & 0.656 & 0.256 & 0.437 & 0.692 \\
Shell & 20 & 0 & 0.663 & 0.246 & 0.433 & 0.672 \\
\bottomrule
\end{tabular}
\end{table}
 
The significance of Theorem~\ref{thm:cocycle-pressure} is that it closes the thermodynamic side of the base theory before any completion-based extension is invoked. This makes the later strict-extension theoremlet mathematically cleaner: the extension theorem is not compensating for a failed or undefined thermodynamic base. It is extending a cocycle theory that is already closed.

Equally important is what this theoremlet does \emph{not} do. It does not yet define the packaged-object theory \(\Tone\), it does not prove any broader-class pressure theorem, and it does not identify the final thermodynamic consequence of the paper. Those belong to the later completion/extension and conditional-disintegration sections.
 \section{Strict theory extension}\label{sec:strict-extension}

We now state the third closed theoremlet in the paper: the strict theory extension theoremlet. This is the point where the paper moves from a closed thermodynamic base theory to a deeper theorem object. The key claim is not that the paper has found a broader theorem class, but that the packaged-object theory \(\Tone\) carries object structure not definable from the cocycle theory \(\Tzero\) on the same audited shell.

\subsection{From the base theory to the extension problem}

Section~\ref{sec:cocycle-pressure} closed the cocycle-pressure theory \(\Tzero\) on the audited shell. Section~\ref{sec:canonical-object} then introduced the packaged-object map
\[
\pi_1:\mathcal{S}_{\mathrm{aud}}\to \mathcal{O}_1
\]
alongside the cocycle-level object map
\[
\pi_0:\mathcal{S}_{\mathrm{aud}}\to \mathcal{O}_0.
\]
The extension problem is therefore intrinsic: does the packaged-object map \(\pi_1\) factor through the cocycle-level object map \(\pi_0\), or does the completion-based theory define additional object structure \cite{LindMarcus1995SymbolicDynamics,Walters1986RelativePressure}?

If there exists a map
\[
\phi:\mathcal{O}_0\to\mathcal{O}_1
\]
such that
\[
\pi_1=\phi\circ \pi_0
\qquad \text{on } \mathcal{S}_{\mathrm{aud}},
\]
then the completion theory adds no new object depth. The packaged-object structure would simply be reconstructible from the base cocycle theory. The theoremlet below states that this does \emph{not} happen on the audited shell.

\subsection{The extension mechanism}

Three ingredients drive the extension.

\paragraph{Saturation.}
For fixed lens \(\ell\) and timescale \(\tau\), the completion endomap
\[
E_{\tau,\ell}(\mu)=U_{\ell}\!\bigl(Q_{\ell}(\mu K^\tau)\bigr)
\]
may saturate: repeated completion stops producing new packaged objects at the current theory depth. Saturation is therefore not the end of the story, but the signal that the current object vocabulary has stabilized.

\paragraph{Material \texorpdfstring{$P4\leftarrow P5$}{P4<-P5} forcing.}
Once saturation occurs, packaging feeds back into the lens layer. Material \(P4\leftarrow P5\) forcing changes the active lens and reveals new packaged strata that were not visible at the previous theory depth. This is the mechanism by which the extension moves from the base theory to the packaged-object theory.

\paragraph{Macro-admissibility obstruction.}
The cocycle-level theory is too coarse to support closed macro-dynamics on the packaged future. This macro-admissibility obstruction is not an accidental diagnostic. It is one of the certificates that the base theory cannot already express the packaged-object structure, and it is conceptually closest to a failure of admissible aggregation or lumpability at the macro level \cite{KemenySnell1976FiniteMarkovChains}.

Together, saturation, forcing, and macro-admissibility obstruction provide the structural route from \(\Tzero\) to \(\Tone\) \cite{TsiokosFoundations}.

\begin{theorem}[Strict theory extension on the audited shell]\label{thm:strict-extension}
Let \(\mathcal{S}_{\mathrm{aud}}\) be the audited shell-stable class, let
\[
\pi_0:\mathcal{S}_{\mathrm{aud}}\to \mathcal{O}_0
\]
be the cocycle-level object map of the base theory \(\Tzero\), and let
\[
\pi_1:\mathcal{S}_{\mathrm{aud}}\to \mathcal{O}_1
\]
be the packaged-object map defined by the completion theory \(\Tone\). Assume:
\begin{enumerate}
    \item the cocycle-pressure theory \(\Tzero\) is closed on \(\mathcal{S}_{\mathrm{aud}}\),
    \item the completion endomap is well-defined on \(\mathcal{S}_{\mathrm{aud}}\),
    \item saturation of the completion dynamics occurs on the audited shell,
    \item material \(P4\leftarrow P5\) forcing produces new packaged strata after saturation,
    \item and the cocycle-level theory exhibits macro-admissibility obstruction on the packaged future.
\end{enumerate}
Then \(\Tone\) is a strict theory extension of \(\Tzero\) on \(\mathcal{S}_{\mathrm{aud}}\). Equivalently, there is no factorization
\[
\pi_1=\phi\circ \pi_0
\qquad \text{on } \mathcal{S}_{\mathrm{aud}}
\]
for any map \(\phi:\mathcal{O}_0\to\mathcal{O}_1\).
\end{theorem}

\begin{proof}[Proof sketch]
The proof follows the closed strict-extension package and is organized around factorization failure rather than family broadening.

\smallskip
\noindent\emph{Step 1: the base theory is already closed.}
By Theorem~\ref{thm:cocycle-pressure}, the cocycle-pressure object of \(\Tzero\) is well-defined on the audited shell. Thus the later extension result is not compensating for a defective or undefined base theory.

\smallskip
\noindent\emph{Step 2: completion defines packaged-object strata.}
The completion endomap \(E_{\tau,\ell}\) is well-defined on the shell and yields persistent packaged strata. These strata define the packaged-object map \(\pi_1\).

\smallskip
\noindent\emph{Step 3: saturation plus forcing produces genuinely new packaged strata.}
Saturation shows that the current completion rule has stabilized at the present theory depth. Material \(P4\leftarrow P5\) forcing then changes the lens and reveals new packaged strata. Thus the extension is not merely an iteration of a fixed completion rule, but a genuine change in theory depth.

\smallskip
\noindent\emph{Step 4: non-definability from \(\Tzero\).}
If \(\pi_1\) factored through \(\pi_0\), then the packaged-object strata revealed after forcing would already be reconstructible from the cocycle-level object map. The audited-shell non-definability step rules this out: shell states with the same \(\Tzero\)-object can carry different packaged-object strata in \(\Tone\).

\smallskip
\noindent\emph{Step 5: macro-admissibility obstruction.}
The failure of macro-admissibility at the cocycle level shows that \(\Tzero\) is too coarse to support closed macro-dynamics on the same packaged future. This obstruction is therefore a positive certificate that the packaged-object theory is not already contained in the base theory.

Combining these points yields the claimed failure of factorization and hence the strict theory extension theoremlet.
\end{proof}

\begin{table}[t]
\caption{Strict theory-extension support on the audited shell.
The packaged-object map fails to factor through the cocycle-level object map
on both witnesses, material $P4\!\leftarrow\!P5$ forcing is present,
and macro-admissibility obstruction is global on the audited panels.
This table supports only the strict theory extension theoremlet
on the audited shell.}
\label{tbl:strict-theory-extension}
\centering
\small
\begin{tabular}{@{}l c c c c c c@{}}
\toprule
Witness & Non-fact.\ rate & $T_1$/$T_0$ strata & Distinct $T_1$ & Forcing & Inadmiss.\ & Sat.\ panels \\
\midrule
Base  & 0.833 & 10 & 30 & 36/36 & 36/36 & 2/9 \\
Shell & 0.833 & 10 &  8 & 36/36 & 36/36 & 8/9 \\
\bottomrule
\end{tabular}
\end{table}
 
Theorem~\ref{thm:strict-extension} is the paper’s open-endedness theoremlet. It shows that the packaged-object theory is not merely a more elaborate visualization of the cocycle-pressure theory. It defines a deeper theory on the same audited shell.

We stress again what the theoremlet does \emph{not} claim. It does not claim a broader theorem class beyond the audited shell. It does not claim a shell-general extension theorem. It does not claim an external/non-SFT breadth theorem. The result is instead a theorem of theory depth: the packaged-object theory \(\Tone\) is strictly richer than the cocycle theory \(\Tzero\) even on a fixed shell-stable class.
 \section{Conditional pressure disintegration}\label{sec:conditional-disintegration}

We now state the fourth closed theoremlet in the paper: the thermodynamic consequence theoremlet. The key point is that the relevant thermodynamic consequence of the strict extension is \emph{not} direct stratumwise root separation. The consequence object used here is a \emph{conditional disintegration} of the base cocycle pressure over the packaged-object fibers of the extended theory \cite{Rokhlin1949FundamentalIdeas}.

\subsection{What is being decomposed}

Section~\ref{sec:cocycle-pressure} closed the cocycle pressure object of the base theory \(\Tzero\). Section~\ref{sec:strict-extension} then showed that the packaged-object theory \(\Tone\) is a strict extension of \(\Tzero\) on the audited shell. The thermodynamic question is therefore not whether different packaged strata behave as independent thermodynamic systems. Rather, the question is whether the cocycle-level pressure is thermodynamically sufficient for the packaged future.

The answer is negative on the audited shell. The packaged object carries thermodynamic information that is not already exhausted by the cocycle-level object map. The correct consequence is therefore a package-conditioned pressure gap.

\subsection{Why the direct stratumwise route is rejected}

A tempting but ultimately incorrect route would be to assign independent pressure roots to each persistent packaged stratum and try to prove direct root separation inside a fixed \(\Tzero\)-object class. That is \emph{not} the route taken here. The reason is conceptual, not merely technical: packaged strata are fiber-averages of the same underlying lawful dynamics and need not be thermodynamically autonomous objects. The paper therefore does not claim a direct stratumwise root-separation theorem.

Instead, the thermodynamic object is global first and conditional second. One starts with the global cocycle pressure \(P_{\Tzero}(s)\), then disintegrates it over the packaged fibers defined by \(\Tone\). Conceptually this is closer to relative pressure and disintegration over factor fibers than to independent-per-stratum pressure roots \cite{LedrappierWalters1977RelativisedVP,Walters1986RelativePressure}.

\subsection{The consequence object}

Let \(\mathcal{F}(x)\) be the family of persistent packaged strata at a shell state \(x\in \mathcal{S}_{\mathrm{aud}}\), and let \(w_{\Sigma}(x)\) be the corresponding package weights. The consequence quantity is
\[
\Delta_{\Tzero\to \Tone}(s)
:=
P_{\Tzero}(s)
-
\sum_{\Sigma\in\mathcal{F}(x)} w_{\Sigma}(x)\,P_{\Sigma}(s),
\]
where \(P_{\Sigma}(s)\) denotes the package-conditioned pressure profile attached to the stratum \(\Sigma\).

The meaning of \(\Delta_{\Tzero\to \Tone}(s)\) is that it measures the thermodynamic insufficiency of the base theory for the packaged future. If \(\Delta_{\Tzero\to \Tone}(s)\) vanished identically, then the cocycle-level pressure would already be sufficient for the packaged thermodynamic future. The theoremlet below states that, on the audited shell, this does not happen.

\subsection{Closure-deficit interpretation}

The pressure gap above is the paper’s thermodynamic analogue of a closure-deficit quantity. It measures how much information about the packaged future is not captured at the cocycle level. In this sense the theoremlet should be read as a conditional consequence of strict theory extension: once the packaged object is strictly richer than the cocycle object, the global cocycle pressure need not remain sufficient for the packaged future.

The section does \emph{not} claim an exact KL-based closure-deficit identity as a closed theorem. Such KL-style quantities remain support-only diagnostics in the present package \cite{CoverThomas2006InfoTheory,TsiokosCastStone}. The closed theoremlet is formulated on the weighted package-conditioned pressure gap itself.

\begin{theorem}[Conditional pressure disintegration on the audited shell]\label{thm:conditional-disintegration}
Let \(\mathcal{S}_{\mathrm{aud}}\) be the audited shell-stable class, let \(P_{\Tzero}(s)\) be the closed cocycle pressure object of the base theory \(\Tzero\), and let \(\mathcal{F}(x)\) be the persistent packaged strata of the extended theory \(\Tone\) with weights \(w_{\Sigma}(x)\). Define
\[
\Delta_{\Tzero\to \Tone}(s)
:=
P_{\Tzero}(s)
-
\sum_{\Sigma\in\mathcal{F}(x)} w_{\Sigma}(x)\,P_{\Sigma}(s).
\]
Assume:
\begin{enumerate}
    \item the cocycle pressure object of \(\Tzero\) is closed on \(\mathcal{S}_{\mathrm{aud}}\),
    \item the packaged-object theory \(\Tone\) is well-defined on \(\mathcal{S}_{\mathrm{aud}}\),
    \item the strict theory-extension theoremlet holds on \(\mathcal{S}_{\mathrm{aud}}\),
    \item and the package-conditioned pressure profiles are well-defined on the audited shell.
\end{enumerate}
Then the weighted package-conditioned pressure gap is nontrivial on the audited shell. Equivalently, the cocycle-level pressure is not sufficient for the packaged future, and the thermodynamic consequence of the extension is a nontrivial conditional pressure disintegration.
\end{theorem}

\begin{proof}[Proof sketch]
The proof is built on the canonical hybrid object and the already-closed theoremlets.

\smallskip
\noindent\emph{Step 1: the global base pressure is closed.}
By Theorem~\ref{thm:cocycle-pressure}, the cocycle pressure \(P_{\Tzero}(s)\) is well-defined on the audited shell.

\smallskip
\noindent\emph{Step 2: package-conditioned profiles are well-defined.}
The packaged-object theory \(\Tone\) supplies the family of persistent packaged strata \(\mathcal{F}(x)\), and the corresponding weighted package-conditioned pressure profiles \(P_{\Sigma}(s)\) are well-defined on the shell.

\smallskip
\noindent\emph{Step 3: strict extension implies thermodynamic insufficiency.}
By Theorem~\ref{thm:strict-extension}, the packaged-object map is not definable from the cocycle-level object map on the audited shell. Thus the packaged future carries structure that is not already exhausted by the base cocycle object.

\smallskip
\noindent\emph{Step 4: the correct consequence is conditional, not independent-per-stratum.}
Because the packaged strata are not treated as autonomous thermodynamic systems, the relevant comparison is not direct root separation across strata. Instead it is the weighted disintegration of the global pressure over the packaging fibers. This is precisely what \(\Delta_{\Tzero\to \Tone}(s)\) measures.

\smallskip
\noindent\emph{Step 5: nontriviality of the gap.}
The audited-shell support package shows that the weighted package-conditioned pressure gap is bounded away from zero on the shell. Hence the cocycle-level pressure is thermodynamically insufficient for the packaged future, and the conditional disintegration consequence is nontrivial.

These points prove the theoremlet.
\end{proof}

\begin{table}[t]
\caption{Conditional pressure disintegration on the audited shell.
The weighted package-conditioned pressure gap is bounded away from zero
on both witnesses, showing that the base cocycle pressure is not sufficient
for the packaged future.
This table supports only the conditional pressure disintegration
theoremlet on the audited shell.}
\label{tbl:conditional-disintegration}
\centering
\small
\begin{tabular}{@{}l c c c c@{}}
\toprule
Witness & Descriptors & Min.\ weighted gap & Gap $>0$ & Inadmiss.\ panels \\
\midrule
Base  & 6 & 0.0735 & yes & 18/18 \\
Shell & 6 & 0.0580 & yes & 18/18 \\
\bottomrule
\end{tabular}
\end{table}
 
Theorem~\ref{thm:conditional-disintegration} is the thermodynamic consequence theoremlet used in the paper. It shows that the strict theory extension is not merely structural. It changes the thermodynamic content of the object theory in a measurable way.

We again emphasize the scope. The theoremlet is proved on the audited shell only. It does not claim a broader-class thermodynamic theorem, a shell-general theorem, an external/non-SFT breadth theorem, or a direct stratumwise root-separation result. The contribution is a closed conditional-disintegration consequence on the canonical hybrid object.
 \section{Related work}\label{sec:related-work}

The mathematical background of this paper lies first in the classical study of Cantor and fractal objects generated by fixed symbolic, self-similar, or graph-directed rules. Standard references include Falconer’s general account of fractal geometry and dimension theory, the original IFS and Moran-construction sources, and the graph-directed and symbolic framework developed by Mauldin and Urba\'nski \cite{Hutchinson1981Fractals,Moran1946AdditiveFunctions,FalconerFractalGeometry,MauldinUrbanskiGDMS}. In those settings, the central object is typically specified in advance by a fixed rule system, and the main questions concern geometry, dimension, pressure, and invariant measures for that fixed object \cite{Bowen1979Quasicircles,Ruelle1978ThermodynamicFormalism,WaltersErgodic1982}.

A second background line is thermodynamic formalism beyond strictly additive observables. The pressure side of the present paper is much closer in spirit to nonadditive or almost-subadditive thermodynamic constructions than to a fixed self-similar pressure formula, even though our final object is not presented in that language alone. For general background on nonadditive thermodynamic formalism, see \cite{BarreiraNonadditiveTF,CaoFengHuang2008SubAdditive,Mummert2006AlmostAdditive,Kingman1968Subadditive}. The present paper does not attempt to prove a broad nonadditive formalism for arbitrary evolving Cantor substrates; instead it closes a pressure object on a specific audited shell-stable class.

A third line of context is the study of local or context-dependent iterated constructions. Recent work on local iterated function systems provides a useful point of comparison because it highlights how admissible compositions may depend on local geometry rather than a single global IFS rule \cite{OliveiraVarandasLocalIFS,BarnsleyDemkoEltonGeronimo1988,FalconerRandomFractals1986,RempeGillenUrbanski2016NonAutonomousCIFS}. However, the theorem package closed in this paper is not a broad local-IFS theorem. We do not claim a shell-general local-IFS theorem or a general non-SFT breadth theorem. Our claims are restricted to the audited shell on which the canonical hybrid object is closed.

Within the prior six-birds program, three earlier papers provide the relevant background and vocabulary. The first is the foundational closure/saturation/extension formalism of \cite{TsiokosFoundations}. The second is the substrate-generation perspective developed in \cite{TsiokosCreateStone}, which motivates the use of a continuous six-mechanism substrate rather than a fixed discrete pipeline. The third is the closure-deficit / conditional-disintegration perspective in \cite{TsiokosCastStone}, which is the closest conceptual antecedent of the thermodynamic consequence theoremlet proved here. These papers are cited only for background and historical context. The present paper is not a general six-birds framework paper; it is a Cantor mathematics paper built on a specific canonical hybrid object.

This distinction is important for novelty. The continuous six-mechanism substrate machinery is adapted from prior six-birds/PICA-style work, and the paper does not claim that substrate-generation mechanism itself as its main novelty. What is new in this branch is the closed theorem package supported by that substrate on the audited shell: a continuous full-loop lawfulness theoremlet, a cocycle pressure closure theoremlet, a strict theory-extension theoremlet, and a conditional pressure disintegration theoremlet on a canonical hybrid Cantor object. The paper’s novelty is therefore a theorem of theory depth rather than broader family coverage.

Finally, we note two routes that are \emph{not} the paper’s claim. First, the paper does not prove a broader-class theorem beyond the audited shell, even though several wider routes were explored during development. Second, the paper does not prove a direct stratumwise root-separation theorem. The thermodynamic consequence is instead formulated, in the paper-native way, as a conditional disintegration of the base pressure over packaging fibers \cite{Rokhlin1949FundamentalIdeas,LedrappierWalters1977RelativisedVP,CoverThomas2006InfoTheory}.
 \section{Discussion and future work}\label{sec:discussion}

The theorem package proved in this paper is intentionally narrow in scope and correspondingly strong in closure. The main result is not a broad theorem about all context-dependent Cantor systems, nor a shell-general theorem, nor a theorem of non-SFT breadth. It is a closed package on an audited shell-stable class for a canonical hybrid Cantor object. The mathematical novelty is therefore a theorem of theory depth rather than a theorem of broad family coverage.

\subsection{What is closed in the present paper}

The paper closes four theoremlets on the audited shell:
\begin{enumerate}
    \item a continuous full-loop lawfulness theoremlet,
    \item a cocycle pressure closure theoremlet,
    \item a strict theory extension theoremlet, and
    \item a conditional pressure disintegration theoremlet.
\end{enumerate}
Taken together, these theoremlets show that the audited shell carries a lawful continuous substrate, that the cocycle-level pressure object is closed there, that the completion-based packaged-object theory is a strict extension of the cocycle theory, and that this extension has a nontrivial thermodynamic consequence.

This goes beyond a simulation-level claim. The paper isolates a canonical theorem object, proves closure of a base pressure theory on that object, proves strict extension of that theory by packaging completion, and proves a conditional pressure disintegration consequence on the same shell-stable class.

\subsection{What the paper does not claim}

The scope restrictions are part of the mathematical content and should be read as theorem-level boundaries, not merely as stylistic caution. In particular, the paper does \emph{not} claim:
\begin{itemize}
    \item a broader-class theorem beyond the audited shell,
    \item a shell-general theorem,
    \item an external or non-SFT breadth theorem beyond what is actually closed here,
    \item a direct stratumwise root-separation theorem,
    \item a packaging-induced broader theorem class claim,
    \item or an exact KL-based closure-deficit theorem.
\end{itemize}
The support-only diagnostics retained in the appendix and internal package do not enlarge these claims. They accompany the theorem package, but they do not substitute for it and they do not widen its scope.

\subsection{Why audited-shell scope is still mathematically meaningful}

One possible reaction to the shell restriction is that the paper proves only a local result around a specially chosen regime. However, the shell restriction is precisely what allows the theorem object to remain mathematically coherent while preserving the full six-mechanism dynamics. A broader class statement that silently drops nondegenerate lens competition, packaging competition, budget activity, or timescale adaptation would address a different and weaker object.

For this reason the paper should be read as establishing a clean theorem package on a canonical object whose scope is audited and explicit.

\subsection{Reserve and stretch routes}

Several mathematically meaningful routes remain outside the closed paper core.

\paragraph{Packaging completion as a broader class route.}
The packaging-completion object is real and mathematically nondegenerate, but the present paper does not claim that it yields a broader theorem class than the cocycle route on the audited shell. That possibility remains a reserve route.

\paragraph{Hybrid-object broadening beyond the audited shell.}
The paper does not claim that the canonical hybrid object already supports a shell-general theory. Whether the extension and disintegration theoremlets can be widened to a larger shell or a more canonical class remains open.

\paragraph{Wider parameter-shell results.}
The present result does not prove that the audited-shell class can be replaced by a substantially wider parameter shell. That remains future work.

\paragraph{External/non-SFT breadth.}
The paper does not claim that the present theorem package already establishes a non-SFT breadth result in the strong external sense. That ambition remains explicitly outside the current theorem package.

\paragraph{Alternative thermodynamic consequence routes.}
The direct stratumwise route was investigated and rejected as the wrong thermodynamic object for the present paper. The accepted route is conditional pressure disintegration. Whether stronger package-conditioned or measure-theoretic consequences can be obtained remains open.

\subsection{Future work}

There are four natural future directions.

\begin{enumerate}
    \item \textbf{Broader shell stability.}
    One can ask whether the audited shell can be widened while preserving all four theoremlets simultaneously.

    \item \textbf{A sharper theorem of non-definability.}
    The strict-extension theoremlet is closed on the audited shell, but a more intrinsic and possibly more canonical non-definability argument would strengthen the overall package.

    \item \textbf{Further thermodynamic consequences.}
    The paper closes conditional pressure disintegration, but stronger measure-, regularity-, or root-level consequences may exist on top of the same canonical object.

    \item \textbf{Broader object classes.}
    The most ambitious direction is to enlarge the theorem package beyond the present audited shell without collapsing the six-mechanism object structure that makes the result interesting in the first place.
\end{enumerate}

\subsection{Final perspective}

The paper’s contribution is therefore best understood as follows. It does not offer the broadest possible theorem class. Instead, it identifies and closes the right theorem object on a lawful shell: a base cocycle theory, a completion-based extension, and a thermodynamic consequence of that extension. If future work succeeds, the natural direction is not to abandon this object in favor of a broader but weaker statement, but to widen the class while preserving the same object-level architecture.
 
\appendix
\section{Supporting evidence}\label{sec:appendix-supporting-evidence}

This appendix records the computational and audit evidence that accompanies the closed theorem package. Its role is supportive rather than foundational. None of the diagnostics below should be read as a substitute for a closed theorem statement, and none of them enlarges the claim scope of the main text.

The supporting evidence used in the paper falls into five categories:
\begin{enumerate}
    \item the continuous pilot,
    \item the regime diagnostics,
    \item the pressure-closure checks,
    \item the strict-extension audit,
    \item and the conditional-disintegration support checks.
\end{enumerate}
All of this evidence is restricted to the audited shell. No shell-general, broader-class, or external/non-SFT breadth claim is inferred from it.

\paragraph{Continuous pilot and regime diagnostics.}
The continuous pilot and the shell-level regime diagnostics establish that the audited-shell class is not a trivial frozen regime. The manuscript-facing summaries of this evidence are Figure~\ref{fig:full-loop-regime} and Figure~\ref{fig:primitive-knockout-closure}. Figure~\ref{fig:full-loop-regime} summarizes persistent full-loop activity on the audited shell, while Figure~\ref{fig:primitive-knockout-closure} summarizes the leave-one-primitive-out degradations that certify six-mechanism causal closure. These figures support the lawfulness theoremlet, but they are not themselves theorem statements.

\paragraph{Pressure-closure checks.}
The finite-horizon shell checks that accompany the cocycle-pressure theoremlet are summarized in Table~\ref{tbl:pressure-closure-support}. Their role is to document the shell-uniform control, bounded Fekete gap proxy, and nondegenerate pressure regime used in the closed pressure package. They support the cocycle pressure closure theoremlet on the audited shell only. They do not constitute a broader-class pressure theorem.

\paragraph{Strict-extension audit.}
The audited-shell extension evidence is summarized in Table~\ref{tbl:strict-theory-extension}. The table records non-factorization, material \(P4\leftarrow P5\) forcing, packaged-strata multiplicity over \(\Tzero\)-classes, and macro-admissibility obstruction. Its purpose is to document the shell-level certificates behind the strict theory extension theoremlet. It does not support any broader-class or shell-general extension claim.

\paragraph{Conditional-disintegration support.}
The shell-level thermodynamic consequence evidence is summarized in Table~\ref{tbl:conditional-disintegration}. The table records the weighted package-conditioned pressure gap on the audited shell and shows that the gap is bounded away from zero on the two frozen witnesses. This supports the conditional pressure disintegration theoremlet. It is important that this support be read in the paper’s chosen language: the theoremlet is a conditional-disintegration statement, not a direct stratumwise root-separation statement.

\paragraph{Support-only closure-deficit proxy.}
The only support-only item retained in the appendix is the closure-deficit proxy table:
\begin{table}[t]
\caption{Support-only closure-deficit proxy values associated with the conditional
disintegration theoremlet.
These KL-style quantities help interpret the packaged future as
thermodynamically richer than the base cocycle object, but they are
not themselves part of any closed theorem statement in this paper.}
\label{tbl:closure-deficit-proxy}
\centering
\small
\begin{tabular}{@{}l c c@{}}
\toprule
Witness & Max.\ closure-deficit proxy & Status \\
\midrule
Base  & $3.77\times 10^{-8}$  & support-only \\
Shell & $1.08\times 10^{-11}$ & support-only \\
\bottomrule
\end{tabular}
\end{table}
 This table is included because it helps interpret the packaged future as thermodynamically richer than the base cocycle object. However, it is \emph{not} part of any closed theorem statement in the paper. In particular, the manuscript does not claim an exact KL-based closure-deficit theorem.

\paragraph{Scope discipline.}
The diagnostics in this appendix are intentionally subordinate to the theoremlets proved in the main text. They support:
\begin{itemize}
    \item the continuous full-loop lawfulness theoremlet,
    \item the cocycle pressure closure theoremlet,
    \item the strict theory extension theoremlet,
    \item and the conditional pressure disintegration theoremlet.
\end{itemize}
They do not support any broader-class theorem, shell-general theorem, external/non-SFT breadth theorem, direct stratumwise root-separation theorem, or packaging-induced broader theorem-class claim. Those remain explicit non-claims of the paper.
 
\bibliographystyle{plainnat}
\bibliography{bib/references}

\end{document}
